\section{Introduction}

Property graph databases have become ubiquitous for modeling complex 
relationships in domains such as financial transaction networks, social 
platforms, knowledge graphs, and fraud detection. Organizations are increasingly 
deploying these systems in cloud environments due to the many benefits of
cloud but cloud deployments by default reveal sensitive data contents to the
cloud provider. While data encryption preserves the confidentiality of 
application data at rest, it removes the cloud provider's ability to process application 
user requests as it cannot compute on encrypted data. Trusted Execution 
Environments (TEEs) have emerged as a way forward to enable a cloud provider
to compute on application data without learning the contents of the data by provideing hardware-level isolation.

However, an adversary, having access to the TEE's execution environment, can 
infer sensitive information by monitoring the program control flow or the 
memory locations accessed during query processing, even without looking at the 
actual data values ~\cite{islam2012access, grubbs2019learning, kellaris2016generic, kornaropoulos2019data, cash2015leakage, demertzis2020seal, grubbs2018pump, gui2019encrypted, lacharite2018improved, oya2021hiding, poddar2020practical, zhang2016all, oya2022ihop}. For 
example, accessing certain subsets of an encrypted transportation network 
frequently might reveal critical routes or hubs, or in an encrypted social 
network, specific patterns of interactions such as frequent queries between two 
specific users might allow an attacker to re-identify those users. This class of 
side-channel attacks has motivated the development of \emph{oblivious} (or
rather \textit{doubly} oblivious) algorithms, where the trace of memory accesses 
and program control flow are independent of the input data. The core challenge 
we address in this paper is: how to efficiently execute property graph
database queries while maintaining oblivious access patterns.


Graph queries inherently involve traversing 
relationships, e.g., finding accounts connected through transactions, 
identifying influence paths in social networks, and so on. These \textit{multi-
hop} queries, i.e., queries traversing a chain on node-edge-node patterns, 
reveal patterns that directly reflect the underlying graph structure. Prior work has demonstrated that access patterns can reveal query predicates, predicate selectivity, and even reconstruct significant portions of the underlying data~\cite{islam2012access, grubbs2019learning, kellaris2016generic, kornaropoulos2019data, cash2015leakage, demertzis2020seal, grubbs2018pump, gui2019encrypted, lacharite2018improved, oya2021hiding, poddar2020practical, zhang2016all, oya2022ihop, abdelraheem2017inference, dautrich2013compromising, giraud2017practical, grubbs2016breaking, kornaropoulos2022leakage, kornaropoulos2020state, kornaropoulos2021response, liu2014search, markatou2019full, pouliot2016shadow}. For regulated industries handling financial or healthcare data, such leakage can violate compliance requirements, even when the data itself remain encrypted.

% Trusted Execution Environments (TEEs) such as Intel SGX/TDX~\cite{hoekstra2013using,intel_tdx} and AMD SEV-SNP~\cite{sev2020strengthening}


Applying techniques from plaintext graph databases as-is fails to preserve
privacy of the graph data and structure. Consider a simple one-hop traversal: 
given a set of source accounts, find all transactions initiated by these 
accounts and their destination accounts. A standard hash join would build a hash 
table on the source accounts and probe it for each edge. However, the number of 
probes per source node reveals node degrees, even when executed within an
enclave. More sophisticated approaches using Oblivious RAM 
(ORAM)~\cite{goldreich1987towards, stefanov2013path} hide individual accesses, 
but add logarithmic overhead per operation, making multi-hop queries using ORAM
impractical at scale. 
% The core difficulty is that graph structure is inherently reflected in access patterns, so hiding it requires fundamentally rethinking how we process the property graph queries. 


Most acyclic property graph database queries can be decomposed into a set of 
relational operators~\cite{yang2025Graphiti}. Hence, we can reuse existing 
oblivious binary join operators~\cite{krastnikov2020efficient,
zheng2017opaque,eskandarian2019oblidb,arasu2013oblivious,mavrogiannakis2025obliviator,reichert2024menhir} to support graph queries. As an example, a multi-hop query 
can be translated into a multi-way relational join operator, executed by a set of 
binary joins. However, this reveals the intermediate output sizes, which can be 
sensitive. On the other hand, a few recent works~\cite{chamani2023graphos, feng2025enabling} support oblivious processing of specific graph algorithms such as 
BFS, DFS, and minimum spanning tree, which differ from the selective, predicate-driven queries typical of graph databases. 

A recent work by Wei and Kerschbaum~\cite{ruidi2025oblivious}, referred to as WK,
computes oblivious multi-way joins while hiding intermediate result sizes, providing 
strong privacy guarantees for acyclic queries. However, we observe that specific
graph databases properties enable optimizations that cannot be exploited as-is in
the WK scheme. Our system uses the WK scheme as a baseline and a building block
to efficiently support property graph queries.


\emph{Our approach.}
We propose \system, an oblivious property graph database, designed to execute 
parallelized graph queries. Graph queries typically tend to have a multi-hop pattern
traversing a sequence of nodes $n1,\cdots,n_i$ connected by edges $e_1,\cdots,
e_{i-1}$. A multi-hop query can be broken down into a series of one-hop queries of 
type $n_i \rightarrow e \rightarrow n_j$. Our key insight is to develop an efficient 
one-hop operator and apply it in parallel where possible within a multi-hop query to 
speed-up processing. 


In developing a one-hop operator, we leverage two properties: First, plaintext 
production
graph databases such as ~\cite{janusgraph, feng2023kuzu, deutsch2019tigergraph, sun2015sqlgraph}, maintain two copies of each edge, in both \textit{forward} and 
\textit{reverse} directions, where the forward edge list is sorted on the source
node ids and the reverse edge list is sorted on the destination node ids. \system
leverages this to process a one-hop query obliviously from both directions, i.e., 
$n_i \rightarrow e$ and $n_j \leftarrow e$, and merge the two results to form the
output on $n_i \rightarrow e \rightarrow n_j$.


Second, we leverage the \textit{offline/online} model wherein the system pre-processes
data in an offline phase to reduce the online processing time, triggered when a
client query arrives. Because graph databases support stored-procedure-like query
templates~\cite{deutsch2019tigergraph, neo4j, janusgraph, bebee2018amazon}, these
systems know a priori the types of nodes and edges to queried, just the query
predicates or the number of hops is provided at query time by the client. \system
utilizes this knowledge in pre-processing the node tables in an offline phase such
that at query time (i.e., online phase) the node and edge attributes are combined
efficiently.


Based on the one-hop operator, we develop a query decomposition framework that 
rewrites multi-hop graph queries to maximize the use of one-hop operators in 
parallel, and uses the multi-way join~\cite{ruidi2025oblivious} on the decomposed 
query graph to combine the various decomposed parts. The decomposed query executes 
significantly faster while providing identical privacy guarantees i.e. all 
intermediate results are either public sizes or hidden. Our evaluation on synthetic 
banking transaction dataset shows 1.5-2x speedup over non-decomposed oblivious execution for 4-hop chain and branch query patterns (see results section).